%% kapa: exam=1 round=21 year=2010 period=1 form=A subject=civil_law subject_ko=민법 seq=1 number=1
%% source: Q-Net 감정평가사 기출문제 (cst00309) artlSeq=5024181 file=extracted/1008/5024181/2010년 제21회 감정평가사 1차시험 기출문제/1교시.hwp
%% conversion: kordoc 4.24.51 -> tools/gen1.py (auto); review=proofread-sonnet; flags=none
\begin{question}{1}
통정허위표시의 무효는 선의의 제3자에게 대항할 수 없는데, 이 경우 제3자에 해당하지 않는 자는? (다툼이 있으면 판례에 의함)
\begin{choices}
  \item 채권의 가장양도에서의 채무자
  \item 가장소비대차에 기한 채권의 양수인
  \item 가장매매의 매수인으로부터 부동산을 매수하여 소유권이전등기를 마친 자
  \item 가장의 저당권설정등기에 기한 저당권 실행에 의해 부동산을 매수한 자
  \item 채무자(파산자)가 가장채권을 보유하고 있다가 파산선고된 경우의 파산관재인
\end{choices}
\end{question}
\answer{1}
